\documentclass[11pt]{article}
\usepackage[margin=1in]{geometry}
\usepackage{amsmath,amssymb,amsfonts,mathtools}
\usepackage{array,booktabs,enumitem,graphicx,xcolor,hyperref}
\usepackage[T1]{fontenc}
\usepackage[utf8]{inputenc}
\title{Picard-Fuchs equations and mirror maps for hypersurfaces}
\author{Source-backed arXiv sample}
\date{}
\begin{document}
\maketitle
\section{Extracted Lists}

The list environments below are extracted from arXiv LaTeX source and wrapped for pdfLaTeX validation.

\begin{enumerate}
\item
If we analytically continue along
 a simple loop around $z=0$ in the counterclockwise direction, $t$
becomes $t+1$.
(It will be convenient to also introduce $q=e^{2\pi it}$, which
remains single-valued near $z=0$.)

\item
There are cycles $\gamma_0$ and $\gamma_1$ such that
$\int_{\gamma_0}\omega(z)$
is single valued near $z=0$, and
\[t=\frac{\int_{\gamma_1}\omega(z)}{\int_{\gamma_0}\omega(z)}\]
in an angular sector near $z=0$.
\end{enumerate}

\begin{enumerate}
\item
There is a period function for $\omega(z)$,
\[f_0(z):=\int_{\gamma_0}\omega(z)\]
which is single-valued near $z=0$.  This period function is unique up to
multiplication by a constant.
(This implies that the cycle $\gamma_0$ is also unique up to a
constant multiple.)

In particular, the family of meromorphic $n$-forms
\[\widetilde\omega(z):=\frac{\omega(z)}{\int_{\gamma_0}\omega(z)}\]
will have the property that
\[\int_{\gamma}\widetilde\omega(z)\equiv1\]
for some $\gamma$, and it is the unique such family up to constant multiple.

\item
Fixing a choice of period function $f_0(z)$ as in part (1), there is a
period function
\[f_1(z):=\int_{\gamma_1}\omega(z)\]
such that $\varphi(z):=f_1(z)/f_0(z)$ transforms as
\[\varphi(z)\mapsto\varphi(z)+1\]
upon transport around $z=0$ in the counterclockwise direction.
The ratio $\varphi(z)$ is unique up to the addition of a constant.
\end{enumerate}
\end{document}
